% Owner: B
We identify the connected components of each graph. A connected component is
a maximal set of nodes in which a path joins every pair of nodes. We measure
distance as the number of edges, or hops, in a shortest path. Following the
course convention, distance is infinite between nodes in different
components \cite{quercini2026}.

We therefore compute diameter and radius only on the largest connected
component, ignoring edge weights. The diameter is the greatest distance
between any pair of nodes in that component. To obtain the radius, we find
each node's greatest distance to another node, then take the smallest of
these values. These measures describe distances within this component only.

An articulation point (cut vertex) is a node whose removal, together with its
incident edges, increases the number of connected components. A bridge (cut
edge) is an edge whose removal increases that number. We find both with
NetworkX \cite{hagberg2008}.
